% =====================================================================
\chapter{Reducción al absurdo y demostración por casos}
\label{tema:4}
% =====================================================================

Algunas afirmaciones no tienen la forma de un condicional, o la tienen pero ni la hipótesis ni la negación de la tesis dan un punto de partida útil. Para ellas se dispone de la reducción al absurdo: suponer lo contrario de lo que se quiere probar y llegar a una contradicción. Este tema la presenta con dos de los teoremas más antiguos de las matemáticas, la irracionalidad de $\sqrt2$ y la infinitud de los números primos, y estudia también la demostración por casos. Se siguen el capítulo~6 y las §4.4–4.5 de~\cite{hammack}, las §1.7–1.8 de~\cite{mcs}, la §3.5 de~\cite{velleman} y el primer \emph{problem set} de CS103~\cite{cs103}.

\section{Reducción al absurdo}

\begin{definicion}[Reducción al absurdo]\label{def:absurdo}
Para demostrar una proposición $P$ por \emph{reducción al absurdo} (o por contradicción) se supone $\neg P$ y se deduce una contradicción, es decir, una proposición de la forma $C\wedge\neg C$. Se concluye $P$.
\end{definicion}

La justificación es que $(\neg P\Rightarrow(C\wedge\neg C))\Rightarrow P$ es una tautología: si $P$ fuera falsa, $\neg P$ sería verdadera, y el condicional $\neg P\Rightarrow(C\wedge\neg C)$, con hipótesis verdadera y tesis falsa, sería falso. Para un enunciado condicional $P\Rightarrow Q$, la negación es $P\wedge\neg Q$ (\cref{prop:leyes}(vi)): se suponen a la vez la hipótesis y la negación de la tesis~\cite[§6.2]{hammack}.

\addtocounter{definicion}{1}%

\addtocounter{definicion}{1}%

\section{La irracionalidad de \texorpdfstring{$\sqrt2$}{raíz de 2}}

Un número real es \emph{racional} si se escribe como $a/b$ con $a,b\in\Z$ y $b\neq 0$, e \emph{irracional} en caso contrario. Todo racional puede escribirse en \emph{forma irreducible}, $a/b$ con $b>0$ y sin divisores comunes mayores que~1: basta dividir $a$ y $b$ por su máximo común divisor. (Alternativamente, se toma la representación con $b$ positivo mínimo, que existe por el principio del buen orden del \cref{tema:5}.)

\begin{teorema}\label{teo:raiz2}
$\sqrt2$ es irracional.
\end{teorema}

La demostración sigue la del teorema~1.8.1 de~\cite{mcs}. Su estructura se repite en los \cref{eje:4-raiz3,eje:4-raiz-cubica,eje:4-log}: suponer una fracción irreducible, obtener una ecuación entre enteros y encontrar un divisor común.

\section{La infinitud de los números primos}

\begin{lema}\label{lem:divisor-primo}
Todo entero $n\ge 2$ tiene al menos un divisor primo.
\end{lema}

\begin{teorema}[Euclides]\label{teo:euclides}
Existen infinitos números primos.
\end{teorema}

\begin{observacion}
La prueba \emph{no} afirma que $N$ sea primo. Con los seis primeros primos, $N=2\cdot 3\cdot 5\cdot 7\cdot 11\cdot 13+1=30\,031=59\cdot 509$ es compuesto; lo que la prueba garantiza es que sus factores primos no están en la lista (\cref{eje:4-euclides-n}).
\end{observacion}

\section{Demostración por casos}

\begin{definicion}[Demostración por casos]\label{def:casos}
Para demostrar $Q$ a partir de $P_1\vee P_2\vee\dots\vee P_k$ se demuestran por separado $P_1\Rightarrow Q$, \dots, $P_k\Rightarrow Q$. Es válido porque $(P_1\vee\dots\vee P_k)\Rightarrow Q\eqlog(P_1\Rightarrow Q)\wedge\dots\wedge(P_k\Rightarrow Q)$, generalización del \cref{eje:1-equivalencias}(c).
\end{definicion}

Lo esencial es que los casos sean \emph{exhaustivos}: su disyunción debe ser verdadera. Las divisiones más habituales son «$n$ par o impar», los restos módulo $m$ y el signo de un número real~\cite[§4.4]{hammack}.

\begin{proposicion}\label{prop:consecutivos}
Para todo $n\in\Z$, el número $n^{2}+n$ es par.
\end{proposicion}

\begin{proposicion}\label{prop:cuadrado-mod4}
Para todo $n\in\Z$, $n^{2}\equiv 0\pmod 4$ o $n^{2}\equiv 1\pmod 4$.
\end{proposicion}

\begin{teorema}[Desigualdad triangular]\label{teo:triangular}
Para todos $x,y\in\R$, $\abs{x+y}\le\abs{x}+\abs{y}$.
\end{teorema}

Cuando dos casos se tratan con el mismo argumento intercambiando los papeles de las variables, se escribe «sin pérdida de generalidad, supóngase $x\le y$» y se trata solo uno~\cite[§4.5]{hammack}. Debe usarse solo cuando la simetría es real.

\addtocounter{definicion}{1}%

\addtocounter{definicion}{1}%

